\subsection{THETA3-C27: a finite cup-product test for the exceptional target and a deformation obstruction}\label{C27}

This is a reduction of the remaining construction problem,
not a construction of an exceptional cycle.  Let $A$ be the
specified fourfold in C25 and $K=\Q(i)$.  Fix a smooth
projective connected surface $X$ over $\overline\Q$ and a
base point.  Put
\[
 M_X=\Hom(\operatorname{Alb}(X),A)\otimes_\Z\Q.
\]
It is a finite-dimensional vector space over $K$, acting
by postcomposition.  Let $m=\dim_K M_X$ and choose a
$K$-basis $f_1,\ldots,f_m$.  Clear denominators so that
these are actual maps $X\to A$ through its Albanese.
Translations will not matter in any cohomological formula.

\paragraph{The exact finite test.}
For one embedding $\sigma:K\to\overline\Q$, choose an
algebraic-Betti basis $v_1,\ldots,v_4$ of
$H^1(A,\overline\Q)_\sigma$ and write
\[
 w_\sigma=v_1\wedge v_2\wedge v_3\wedge v_4.
\]
This spans one Weil line.  The individual $v_i$ are not
asserted to be of pure Hodge type over $\overline\Q$;
the exterior product has type $(2,2)$ because $A$ has
signature $(2,2)$.
Define the homogeneous quartic
\begin{equation}\label{C27-quartic}
 F_X(z_1,\ldots,z_m)
   =\int_X\bigwedge_{r=1}^4
                 \left(\sum_{j=1}^m z_j f_j^*v_r\right)
             \in\overline\Q[z_1,\ldots,z_m].
\end{equation}
Then the following are equivalent:
\begin{enumerate}
\item $F_X$ is not identically zero;
\item some algebraic map $f:X\to A$ has
      $e_W\operatorname{cl}(f_*[X])\ne0$, where $e_W$
      is the explicit algebraic Weil projector of C13.
\end{enumerate}
Thus at most $\binom{m+3}{4}$ explicitly symmetrized
cup-product coefficients decide the existence of such
a map on the given surface.  An individual unsymmetrized
four-linear coefficient is not a sufficient test.
There is also a finite constructive search: if the
quartic is nonzero, one of the at most $5^m$ actual maps
\[
                \sum_j n_jf_j,\qquad
                      n_j\in\{0,1,2,3,4\},
\]
has nonzero Weil projection.  Thus once the isotypic
basis is known, no unbounded denominator search is
needed to select the map.

\textbf{Proof.}
For $a_j\in K$, the rational homomorphism
$f_a=\sum_j a_jf_j$ satisfies
\[
 f_a^*v_r=\sum_j\sigma(a_j)f_j^*v_r,
 \qquad
       \int_X f_a^*w_\sigma
                  =F_X(\sigma(a_1),\ldots,\sigma(a_m)).
\]
If the quartic vanishes on $\{0,1,2,3,4\}^m$, view it
as a polynomial in its last variable.  Its degree in
that variable is at most four, so its values at the
five grid points force all its coefficient polynomials
to vanish on the remaining grid.  Induction on $m$
forces every coefficient to be zero.  Hence a nonzero
quartic has a nonzero grid evaluation.  Since the
chosen $f_j$ are already actual homomorphisms, the map
$f_n=\sum_jn_jf_j$ at this evaluation is actual too.
The projection formula identifies it with the pairing
of $\operatorname{cl}((f_n)_*[X])$ and $w_\sigma$.

In $\bigwedge^4H^1(A,\overline\Q)$ the class
$w_\sigma$ pairs only with the opposite pure Weil line:
wedging with any mixed eigenspace gives more than four
vectors in the four-dimensional $\sigma$ eigenspace.
The pairing with that opposite line is nondegenerate.
The cycle class is rational, so its two pure Weil
components are conjugate.  Hence this pairing is
nonzero exactly when its Weil projection is nonzero.
Conversely every map $X\to A$, after a translation,
comes from its Albanese and therefore lies in $M_X$;
the displayed evaluation gives the converse implication.
If $m=0$, both conditions are false, with the quartic
understood as zero.  The coefficient count is the
number of monomials of degree four in $m$ variables.
\hfill$\square$

When this test succeeds, it gives the actual cycle
$e_W f_*[X]$, and C13 then spans the entire rational
Weil space by algebraic endomorphisms.  It requires
no arbitrary Hodge projector.  Degree-one Hodge maps
used to determine $M_X$ are precisely homomorphisms
of abelian varieties after denominators are cleared;
the unproved datum is the particular nonzero quartic.

\paragraph{What finite covers can and cannot change.}
If $\pi:X'\to X$ is finite \'etale of degree $d$, use
the pulled-back maps $f_j\pi$ on $X'$.  For this specified
subspace of maps one has the exact identity
\begin{equation}\label{C27-cover}
                         F_{X',\pi^*M_X}=dF_X.
\end{equation}
Indeed every integrand is a pullback of the corresponding
top-degree form on $X$, and integration multiplies by
$d$.  In particular, taking another finite cover does
not repair a zero Weil pairing of existing maps.
If $M_{X'}=\pi^*M_X$, no map on $X'$ can repair it.
If the larger space contains new maps, its full quartic
has additional mixed coefficients which must actually
be tested; an increase in multiplicity alone is not
a nonvanishing proof.

For the known C25 Albanese quotient the map factors
through $H$.  The common-cover construction using two
curve projections factors through a product of two
curves.  All maps in the corresponding subspace have
zero quartic by C14.  This assertion concerns that
known subspace, not all of $M_X$ on an uncomputed
cover.  Thus a productive search in the noncongruence
tower must produce genuinely new $A$-isotypic maps
and a nonzero mixed coefficient in
\eqref{C27-quartic}, or use a codimension-two
correspondence outside this map construction.

\paragraph{A fixed ball-source map cannot follow a nontrivial first-order target deformation.}
There is a precise obstruction to one proposed
universal-family approach.  Let $f:X\to A$ generate
$A$, meaning that the induced Albanese homomorphism
is surjective.  Let $A_\epsilon$ be a first-order
deformation of $A$ over $\C[\epsilon]/(\epsilon^2)$.
If $f$ lifts to a map from the constant deformation
$X_\epsilon=X\times\operatorname{Spec}\C[\epsilon]$
to $A_\epsilon$, then
\begin{equation}\label{C27-KS}
                  \operatorname{KS}(A_\epsilon)=0
                           \quad\hbox{in }H^1(A,T_A).
\end{equation}

Here is the cotangent and obstruction calculation.
The exact tangent sequence for deformations of a
morphism is Stacks Project, Lemma~93.10.3,
Section~93.10, Tag~0E3S.  Its last term is
\[
 \operatorname{Ext}^1_X(Lf^*L_{A/\C},\mathcal O_X).
\]
Since $A$ is smooth, $L_{A/\C}=\Omega^1_A$ is locally
free in degree zero.  This term is therefore
$H^1(X,f^*T_A)$.  With the source deformation zero,
the obstruction is the pullback of the target
Kodaira--Spencer class, up to the irrelevant overall
sign in the difference convention.  This identification
can also be verified directly: choose locally trivial
deformations on affine charts of $A$; their transition
maps are $1+\epsilon v_{ab}$.  Local lifts of $f$
differ by sections of $f^*T_A$, and their coboundary
must equal the pulled-back cocycle $f^*v_{ab}$.
Thus a lift forces that cocycle class to vanish.

Invariant vector fields trivialize the tangent bundle,
so this obstruction map is exactly
\[
 \operatorname{Lie}(A)\otimes H^1(A,\mathcal O_A)
  \xrightarrow{\ 1\otimes f^*\ }
 \operatorname{Lie}(A)\otimes H^1(X,\mathcal O_X).
\]
The map $f^*$ on $H^1(\mathcal O)$ is injective:
duality identifies it with the differential of the
dual homomorphism to the surjection
$\operatorname{Alb}(X)\to A$, followed by the
Albanese identification of degree-one cohomology.
Equivalently, a polarization splits that surjection
up to isogeny, giving a left inverse in characteristic
zero.  The obstruction map is therefore injective,
which proves \eqref{C27-KS}.

This applies in particular while retaining a fixed
finite \'etale cover of Schoen's $Y$.  A finite
\'etale cover of the constant first-order deformation
of $Y$ is itself constant with its specified special
fibre, by invariance of finite \'etale covers under
nilpotent thickenings; an exact reference is Stacks
Project, Lemma~58.8.3, Tag~0BQB.
The same obstruction applies to a degree-one
polarized correspondence realizing a surjective
Albanese quotient, because it is such a rational
homomorphism after clearing denominators.

Consequently an attempted deformation from a CM
ball-derived map to the non-CM target cannot retain
that fixed source map in a nonzero target tangent
direction.  A deformation route would have to
construct a genuine relative codimension-two cycle,
with a separate flatness and class-transport argument,
rather than deform this degree-one quotient.  This
does not rule out such higher correspondences or
singular-cycle deformations, and it does not produce
one.  The exact outstanding data are now either
a nonzero finite quartic \eqref{C27-quartic} on a
new isotypic source, or that higher relative cycle.
\hfill$\square$
